% GENERATED FILE. Edit canonical lessons, then run npm run generate:books.
% canonical-source-hash: fnv1a64:01283679eecf83a9
% canonical-lessons: ES-C563-fabula, ES-C563-relato, ES-C563-narrar, ES-C563-relatar, ES-C563-narrador

\chapter{Contar una historia}
\label{ch:es-contar-una-historia}
\hlchaptermodality{563}

\begin{chapteropening}
\textbf{By the end of this chapter:} \emph{I can talk about telling stories: a fable, an account, the narrator, and the verbs to narrate and to relate.}

Complete the last lesson of chapter 563: I can talk about telling stories: a fable, an account, the narrator, and the verbs to narrate and to relate.
\end{chapteropening}

\section[la fábula]{la fábula — the fable}

\label{lesson:ES-C563-fabula}



\begin{quote}
Before the new one: say the Spanish for to burst out, then the Spanish for to move away.
\end{quote}

\subsection*{You'll want to know: la fábula}
\textbf{la fábula} — \textquotedblleft{}the fable\textquotedblright{}.

\begin{quote}
En la fábula, la tortuga le gana a la liebre.
\end{quote}

\begin{quote}
\emph{In the fable, the tortoise beats the hare.}
\end{quote}

\emph{La moraleja} is the moral at the end.

\begin{cousinweb}
\emph{fábula} is a learned borrowing from Latin \emph{fābula} 'a story, a tale' (EtymDB), built on \emph{fārī} 'to speak' (Lewis \& Short). The inherited form of the same word is \emph{habla}.
\end{cousinweb}

\begin{grammarlens}[title={What you've built}]
One more word for telling, explaining and describing at B1.
\end{grammarlens}

\subsection*{Your turn}
\begin{itemize}
\raggedright
  \item \emph{Say it:} \emph{la fábula}
  \item \emph{Say it:} \emph{la fábula}, once more
  \item \emph{Say it:} a sentence of your own with \emph{la fábula}
  \item \emph{Recall:} say the Spanish for to burst out, then the Spanish for to move away, then say \emph{la fábula} again
\end{itemize}

\begin{tcolorbox}[breakable,colback=teal!4,colframe=teal!35!black,title={Before you move on}]
What does la fábula mean? (\textquotedblleft{}The fable\textquotedblright{}.) Say it once more.
\end{tcolorbox}
\section[el relato]{el relato — the account; the short story}

\label{lesson:ES-C563-relato}



\begin{quote}
Before the new one: say the Spanish for to move away, then the Spanish for the fable.
\end{quote}

\subsection*{You'll want to know: el relato}
\textbf{el relato} — \textquotedblleft{}the account; the short story\textquotedblright{}.

\begin{quote}
Escribió un relato sobre la guerra.
\end{quote}

\begin{quote}
\emph{She wrote a short story about the war.}
\end{quote}

\begin{cousinweb}
\emph{relato} comes from Latin \emph{relātus} 'brought back, reported' (EtymDB), the participle of \emph{referre} 'to carry back', built on \emph{ferre} 'to carry'.
\end{cousinweb}

\begin{grammarlens}[title={What you've built}]
One more word for telling, explaining and describing at B1.
\end{grammarlens}

\subsection*{Your turn}
\begin{itemize}
\raggedright
  \item \emph{Say it:} \emph{el relato}
  \item \emph{Say it:} \emph{el relato}, once more
  \item \emph{Say it:} a sentence of your own with \emph{el relato}
  \item \emph{Recall:} say the Spanish for to move away, then the Spanish for the fable, then say \emph{el relato} again
\end{itemize}

\begin{tcolorbox}[breakable,colback=teal!4,colframe=teal!35!black,title={Before you move on}]
What does el relato mean? (\textquotedblleft{}The account; the short story\textquotedblright{}.) Say it once more.
\end{tcolorbox}
\section[narrar]{narrar — to narrate, to tell}

\label{lesson:ES-C563-narrar}



\begin{quote}
Before the new one: say the Spanish for the fable, then the Spanish for the account.
\end{quote}

\subsection*{You'll want to know: narrar}
\textbf{narrar} — \textquotedblleft{}to narrate, to tell\textquotedblright{}.

\begin{quote}
El libro narra la vida de un pescador.
\end{quote}

\begin{quote}
\emph{The book tells the life of a fisherman.}
\end{quote}

A written word; in speech, \emph{contar}.

\begin{cousinweb}
\emph{narrar} is a learned borrowing from Latin \emph{nārrāre} 'to tell, to relate' (EtymDB).
\end{cousinweb}

\begin{grammarlens}[title={What you've built}]
One more word for telling, explaining and describing at B1.
\end{grammarlens}

\subsection*{Your turn}
\begin{itemize}
\raggedright
  \item \emph{Say it:} \emph{narrar}
  \item \emph{Say it:} \emph{narrar}, once more
  \item \emph{Say it:} a sentence of your own with \emph{narrar}
  \item \emph{Recall:} say the Spanish for the fable, then the Spanish for the account, then say \emph{narrar} again
\end{itemize}

\begin{tcolorbox}[breakable,colback=teal!4,colframe=teal!35!black,title={Before you move on}]
What does narrar mean? (\textquotedblleft{}To narrate, to tell\textquotedblright{}.) Say it once more.
\end{tcolorbox}
\section[relatar]{relatar — to relate, to recount}

\label{lesson:ES-C563-relatar}



\begin{quote}
Before the new one: say the Spanish for the account, then the Spanish for to narrate.
\end{quote}

\subsection*{You'll want to know: relatar}
\textbf{relatar} — \textquotedblleft{}to relate, to recount\textquotedblright{}.

\begin{quote}
El testigo relató lo que vio.
\end{quote}

\begin{quote}
\emph{The witness recounted what he saw.}
\end{quote}

\begin{cousinweb}
\emph{relatar} is built on \emph{relato} 'an account', which comes from Latin \emph{relātus} 'reported', the participle of \emph{referre} 'to carry back' (Lewis \& Short).
\end{cousinweb}

\begin{grammarlens}[title={What you've built}]
One more word for telling, explaining and describing at B1.
\end{grammarlens}

\subsection*{Your turn}
\begin{itemize}
\raggedright
  \item \emph{Say it:} \emph{relatar}
  \item \emph{Say it:} \emph{relatar}, once more
  \item \emph{Say it:} a sentence of your own with \emph{relatar}
  \item \emph{Recall:} say the Spanish for the account, then the Spanish for to narrate, then say \emph{relatar} again
\end{itemize}

\begin{tcolorbox}[breakable,colback=teal!4,colframe=teal!35!black,title={Before you move on}]
What does relatar mean? (\textquotedblleft{}To relate, to recount\textquotedblright{}.) Say it once more.
\end{tcolorbox}
\section[el narrador]{el narrador — the narrator}

\label{lesson:ES-C563-narrador}



\begin{quote}
Before the new one: say the Spanish for to narrate, then the Spanish for to relate.
\end{quote}

\subsection*{You'll want to know: el narrador}
\textbf{el narrador} — \textquotedblleft{}the narrator\textquotedblright{}.

\begin{quote}
El narrador es un niño de diez años.
\end{quote}

\begin{quote}
\emph{The narrator is a ten-year-old boy.}
\end{quote}

A woman is \emph{la narradora}.

\begin{cousinweb}
\emph{narrador} comes from Latin \emph{nārrātor} 'a storyteller' (EtymDB), from \emph{nārrāre} 'to tell'.
\end{cousinweb}

\begin{grammarlens}[title={What you've built}]
One more word for telling, explaining and describing at B1.
\end{grammarlens}

\subsection*{Your turn}
\begin{itemize}
\raggedright
  \item \emph{Say it:} \emph{el narrador}
  \item \emph{Say it:} \emph{el narrador}, once more
  \item \emph{Say it:} a sentence of your own with \emph{el narrador}
  \item \emph{Recall:} say the Spanish for to narrate, then the Spanish for to relate, then say \emph{el narrador} again
\end{itemize}

\begin{tcolorbox}[breakable,colback=teal!4,colframe=teal!35!black,title={Before you move on}]
What does el narrador mean? (\textquotedblleft{}The narrator\textquotedblright{}.) Say it once more.
\end{tcolorbox}
